\section{A physical period family and the complete periodic annihilator}
\label{sec:joint-arithmetic}
The four records in Proposition~\ref{prop:lattice} settle the integer
coordinates at zero roof frequency. We now add a family of actual periodic
trajectories that settles every roof frequency. A slightly larger section is
specified explicitly; the original section and every assertion about it
remain as stated in the preceding sections.

\subsection{The enlarged section}
Put
\begin{equation}\label{eq:enlarged-section}
 U=\{(\alpha,p):|\alpha-\pi/6|<3/50,\ |p|<3/20\},
 \qquad Y_R^*=Y_R\cup U,\qquad c_*:=\nu(Y_R^*)=\frac{91}{10000\pi}.
\end{equation}
The rectangle $U$ is disjoint from the four squares defining $Y_R$.
At the square with the same angular center the $p$ coordinate is near
$-1/2$; the other three squares are separated in angle or in $p$.
The area of $U$ is $4(3/50)(3/20)$, and the section density is
$1/(4\pi)$, proving the value of $c_*$. Write
$F_R^*=T_R^{r_R^*}$ for the first return to $Y_R^*$, and let
$\kappa_R^*$ and $\tau_R^*$ be the sums of the physical displacement
and flight length before that return. Its invariant probability is
$\nu_R^*=\nu|_{Y_R^*}/c_*$. An assertion about $F_R$ does not become
an assertion about $F_R^*$ merely by changing notation.

Rotate coordinates through $-\pi/6$ and put
\begin{gather*}
 D=\sqrt3,\qquad O=(0,0),\quad C=(D,0),\quad
 A=(D/2,-1/2),\quad B=(D/2,1/2),\\
 d=\frac12-R,\qquad g=D-2R.
\end{gather*}
The rotated lattice consists of $(Dk/2,l/2)$ with integers $k,l$ of
the same parity. The normal orbit between $O$ and $C$ is unobstructed:
the two intermediate centers $A,B$ have distance $1/2$ from its
supporting line, and all other centers are farther away. It has two
physical flights of total length $2g$, zero winding, and exactly one
visit to $Y_R^*$ per period. Its selected state is $(\pi/6,0)$ in the
unrotated coordinates. The four cycles of Proposition~\ref{prop:lattice}
still have exactly one visit per period. Indeed, the equilateral cycle
has $p=-1/2$, the two-cycle has angles $0,\pi$, and the winding cycles
have $|p|=\sin\theta_R>0.79$. None gains a visit to $U$.

\subsection{A convex optical action}
For $-d\le y,z\le0$ define
\begin{align}
 x_R(y)&=\sqrt{R^2-y^2},\qquad u_R(y)=R-x_R(y),\nonumber\\
 e_R(y)&=\sqrt{(D/2-x_R(y))^2+(d+y)^2},\label{eq:excursion-lengths}\\
 \ell_R(y,z)&=\sqrt{(D-x_R(y)-x_R(z))^2+(z-y)^2}.\nonumber
\end{align}
The function $e_R$ is the length from the top contact
$q_A=(D/2,-d)$ of the disk at $A$ to a contact at height $y$ on the
facing arc of the disk at $O$ or $C$. The function $\ell_R$ is a
flight length between the two facing arcs. For $m\ge1$ let
\begin{equation}\label{eq:excursion-action}
 \mathcal A_{m,R}(y_1,\ldots,y_{2m})
 =e_R(y_1)+\sum_{j=1}^{2m-1}\ell_R(y_j,y_{j+1})+e_R(y_{2m}),
 \qquad L_m(R)=\min_{[-d,0]^{2m}}\mathcal A_{m,R}.
\end{equation}

\begin{lemma}[A uniformly regular physical family]\label{lem:excursion-family}
For each $m\ge1$ and $R\in I$, the minimum in
\eqref{eq:excursion-action} is unique. Its coordinates satisfy
\begin{equation}\label{eq:excursion-height}
 -\frac{2d}{5}<y_j<0,\qquad y_j=y_{2m+1-j}.
\end{equation}
The contacts with these heights form a genuine specular periodic orbit
with center word $A,(O,C)^m,A$. Every flight avoids all third disks.
The orbit has zero winding, $2m+1$ physical collisions, and exactly
$m$ returns to $Y_R^*$. The period length $L_m(R)$ is analytic in $R$
for each fixed $m$.
\end{lemma}
\begin{proof}
We supply the variational, mechanical and return checks separately.
On the whole box, $x_R(y)>11/25$, $x_R(y)\le47/100$, and
$1.73<D<1.74$. The positive horizontal quantities
\[
 a(y)=D/2-x_R(y),\qquad h(y,z)=D-x_R(y)-x_R(z)
\]
satisfy $a>79/200$ and $h>79/100$. The second derivative of
$-x_R(y)$ is $R^2/x_R(y)^3>0$. The Hessian formula for a Euclidean
norm shows that the Hessian of $e_R$ or $\ell_R$ is the sum of a
positive semidefinite term and the Hessian of its horizontal
component multiplied by its positive horizontal direction cosine.
The latter cosine exceeds $9/10$ on this box. Thus each segment
length is strictly convex in its variables; in fact the Hessian of
$\mathcal A_{m,R}$ dominates the identity. One may check the last
numerical bound from
$(9/10)(9/20)^2/(47/100)^3>1$. Compactness and strict convexity give
a unique minimizer.

No coordinate equals $-d$. For such a coordinate $y$ and an incident
interior segment, the derivative is
\[
 \partial_y\ell_R(y,z)=\frac{h(y,z)y/x_R(y)+y-z}{\ell_R(y,z)}<0,
\]
since $z\ge y=-d$. For an end segment,
$e_R'(-d)=-a(-d)d/(x_R(-d)e_R(-d))<0$.
Increasing that coordinate therefore decreases the action.
At an end coordinate equal to zero, $e_R'(0)=d/e_R(0)>0$, and
all other incident derivatives are nonnegative. At an interior zero
adjacent to a negative coordinate, its derivative is positive.
If a block of zero coordinates had no such neighbor it would reach
an endpoint. A variation into the box excludes all these cases.
Every coordinate is consequently in $(-d,0)$. Reflection in the
line $x=D/2$ reverses the center word and preserves the action, so
uniqueness gives the symmetry in \eqref{eq:excursion-height}.

Here is a uniform improvement of the height bound. A smallest
coordinate $y=-w$ cannot be interior, because both incident
derivatives there are negative. Take it to be $y_1$, using reversal
if needed, and write $z=y_2$. Stationarity gives
\[
 d-w=\left(a+e_R(y)\frac{h}{\ell_R(y,z)}\right)\frac{w}{x_R(y)}
           +\frac{e_R(y)}{\ell_R(y,z)}(w+z).
\]
The last term is nonnegative since $z\ge-w$. Moreover,
\[
 \frac{a+e_R(y)h/\ell_R(y,z)}{x_R(y)}
 >\frac{(19/10)(79/200)}{47/100}=\frac{1501}{940}>\frac32.
\]
It follows that $w<2d/5$. This bounds every coordinate.

Interior stationarity of the optical length is precisely equality
of tangential incoming and outgoing velocities at each contact.
The outgoing directions point into the exterior normal half-plane.
For example, at a contact on $O$, its outward normal is
$(x_R(y),y)/R$. Its dot product with a direction to $C$ has numerator
at least $x_R(y)h-|y||z-y|>0$; for a direction to $q_A$ both the
horizontal contribution and $y(-d-y)$ are positive. The checks at
$C$ are reflections of these. The symmetry of the two end contacts
implies the reflection law at $q_A$, whose normal is vertical. There
$e_R(y_1)<11/25$ and $d+y_1>3d/5\ge9/500$, so the normal
incidence cosine exceeds $9/220>1/25$. The other contacts are
uniformly transverse by the preceding horizontal-component bounds. Thus the stationary polygon
is specular rather than the straight-through solution of the
stationarity equation.

The long flights lie in $-2d/5<y<0$, so their distance from the disk
at $A$ is at least $3d/5\ge9/500$, and from the disk at $B$ at
least $d\ge3/100$. End flights join $q_A$ to the appropriate facing
arc and leave both endpoint disks in their exterior supporting
half-planes. Their third disk $O$ or $C$ has horizontal separation
at least $D/2>R$, and $B$ is vertically separated. All segments
lie in the rectangle
\[
       11/25<x<D-11/25,\qquad -d\le y\le0.
\]
For the remaining lattice centers, the possibilities $k=0,2$ with
nonzero $l$ have vertical separation at least $1-d\ge0.95$;
$k=1$, $|l|\ge3$ have separation at least $3/2-d$; and
$k\notin\{0,1,2\}$ have horizontal separation exceeding $1.3$.
This enumerates the whole lattice and proves nonocclusion, not just
a check of a finite list of nearby obstacles.

At an outgoing contact on $O$, the normal angle differs from
$\pi/6$ by at most $|y|/x_R(y)<1/22<3/50$. The outgoing velocity
points toward $C$ and its angular deviation from the horizontal is
at most $|z-y|/h<2/79$. Thus
\[
 |p|\le \frac1{22}+\frac2{79}<\frac3{20}.
\]
All $m$ contacts on $O$ belong to $U$. Contacts on $C$ have angles
near $7\pi/6$, and $q_A$ has angle $2\pi/3$; none belongs to
$Y_R^*$. The return counts are $2$ between successive contacts on
$O$, with one count $3$ through $C,A,O$. Their sum is $2m+1$,
and there are $m$ induced returns. The orbit is closed in the plane,
so its winding is zero. All selected contacts and unselected contacts
have strict membership margins, giving actual regular return branches.
Finally the positive Hessian and the interior minimum allow the
analytic implicit function theorem. This proves analyticity for each
fixed $m$; no bound for high parameter derivatives uniform in $m$ is
being inferred.
\end{proof}

\subsection{Strictly increasing bounded excess lengths}
The following argument supplies an infinite sequence of period relations.
It does not infer nonarithmeticity from finitely many decimal lengths.

\begin{theorem}[Strict period excess]\label{thm:period-excess}
For the family of Lemma~\ref{lem:excursion-family}, define
$E_m(R)=L_m(R)-2m g$. Then, for every $R\in I$,
\begin{equation}\label{eq:period-excess}
 0<E_m(R)<E_{m+1}(R)\le \frac{2d^2}{g}<\frac1{100}.
\end{equation}
In particular $E_{m+1}(R)-E_m(R)>0$ and these differences tend to
zero. The convergence of $E_m$ is used pointwise in $R$; the subsequent
compact-frequency assertion supplies its own uniformity argument.
\end{theorem}
\begin{proof}
Suppress $R$ and set $K(y,z)=\ell(y,z)-g$ on $[-d,0]^2$.
The horizontal component of the flight gives
\begin{equation}\label{eq:kernel-coercivity}
 K(y,z)\ge u(y)+u(z),
\end{equation}
with equality exactly when $y=z$. Define continuous value functions
\[
 V_0(y)=u(y),\qquad
 V_{j+1}(y)=\min_{z\in[-d,0]}\{K(y,z)+V_j(z)\}.
\]
They are nonnegative and vanish at zero. For $j\ge1$, choosing
$z=0$ gives $V_j(y)\le K(y,0)=O(y^2)$ near zero; $V_0$ has the
same bound. The implied bound is independent of $j$.
Equation~\eqref{eq:kernel-coercivity} implies $V_1(y)>V_0(y)$ for
$y<0$: equality at a minimizing $z$ would require both $z=0$ and
equality in \eqref{eq:kernel-coercivity}, hence $y=0$.

For any $y<0$, zero is not a minimizer of $K(y,z)+V_j(z)$. Indeed,
\[
 \partial_zK(y,0)=\frac{-y}{\ell(y,0)}>0,
\]
whereas $0\le V_j(z)=O(z^2)$. A sufficiently small negative $z$
therefore reduces the value. If $V_j>V_{j-1}$ on $[-d,0)$, evaluate
at a minimizer $z_*<0$ for $V_{j+1}(y)$. It follows that
\[
 V_{j+1}(y)=K(y,z_*)+V_j(z_*)
   >K(y,z_*)+V_{j-1}(z_*)\ge V_j(y).
\]
Induction proves the strict inequality for every $j$ and $y<0$.

The minimizing path is symmetric, and its middle edge has length
$\ell(y_m,y_m)=g+2u(y_m)$. Splitting the action at that edge gives
the exact identity
\begin{equation}\label{eq:bellman-excess}
 E_m=2\min_{y\in[-d,0]}\{e(y)-g/2+V_{m-1}(y)\}.
\end{equation}
The outer minimizer cannot be zero, since $e'(0)=d/e(0)>0$ and
$V_{m-1}(y)=O(y^2)$. Evaluating the formula for $E_{m+1}$ at its
negative minimizer, and using strict increase of $V_j$ there, proves
$E_{m+1}>E_m$. Since $e(y)\ge D/2-x_R(y)\ge g/2$, equality
$E_m=0$ would force $y=0$ and then $d=0$, which is excluded.
Taking all heights zero gives
\[
 E_m\le2\sqrt{(g/2)^2+d^2}-g
      =\frac{2d^2}{\sqrt{(g/2)^2+d^2}+g/2}\le\frac{2d^2}{g}.
\]
Here $d\le1/20$ and $g>79/100$, proving the last bound in
\eqref{eq:period-excess}. A bounded increasing sequence has summable
positive increments, which must tend to zero.
\end{proof}

\subsection{Joint arithmetic and its exact scope}
A regular periodic return orbit has a total record $(K,N,T,h)$:
winding $K\in\Z^2$, physical collision count $N$, flight length $T$,
and induced period $h$. Define its phase at
$(u,s,\beta,c)\in\T^2\times\T\times\R\times\T$ to be
$u\cdot K+sN+\beta T-ch$, with $\T=\R/(2\pi\Z)$.

\begin{theorem}[Trivial joint periodic annihilator]\label{thm:joint-periodic}
For every $R\in I$, if
\begin{equation}\label{eq:joint-periodic-phase}
 \exp i(u\cdot K+sN+\beta T-ch)=1
\end{equation}
for every regular periodic orbit of $F_R^*$, then
$\beta=0$ and $u=s=c=0$ on their tori. It suffices to test the four
cycles of Proposition~\ref{prop:lattice}, the long normal two-cycle,
and the family of Lemma~\ref{lem:excursion-family}.
\end{theorem}
\begin{proof}
The long normal cycle gives $\exp i(2s+2\beta g-c)=1$.
The $m$th excursion gives
$\exp i((2m+1)s+\beta L_m-mc)=1$. Dividing the latter by the
$m$th power of the former yields
\[
                    \exp i(s+\beta E_m)=1.
\]
Taking consecutive quotients shows that
$\beta(E_{m+1}-E_m)\in2\pi\Z$ for all $m$. If $\beta\ne0$,
Theorem~\ref{thm:period-excess} makes these quantities nonzero and
strictly smaller than $2\pi$ in absolute value for large $m$, a
contradiction. Hence $\beta=0$. The four unchanged first-return
records in \eqref{eq:matrix} now give $u=s=c=0$.
\end{proof}

\begin{corollary}[Finite uniform separation on compact frequency sets]
\label{cor:compact-periodic}
For every $B<\infty$ and $\varepsilon>0$, there are a finite list of
these physical periodic words and $\eta>0$ such that, uniformly in
$R\in I$ and in $|\beta|\le B$ at distance at least $\varepsilon$
from the trivial phase, at least one word has
$|\exp i(u\cdot K+sN+\beta T-ch)-1|\ge\eta$.
\end{corollary}
\begin{proof}
Each word has fixed integer record and continuous length on $I$.
Theorem~\ref{thm:joint-periodic} gives a nonzero discrepancy at every
point of the compact parameter--frequency set in question. The open
sets on which a given word has positive discrepancy cover that set.
Take a finite subcover and minimize the maximum of its discrepancies.
This minimum is positive and supplies $\eta$.
\end{proof}

\begin{corollary}[Full real period rank]\label{cor:real-period-rank}
The real vectors $(K_1,K_2,N,h,T)$ of five fixed cycles span $\R^5$.
The absolute determinant is $2(\sqrt3-1)$, independently of $R$.
Consequently no nonzero real linear combination of the four induced
record coordinates is cohomologous to a constant by a transfer
function for which periodic evaluation is valid.
\end{corollary}
\begin{proof}
Append lengths $L_2,L_3,L_4,L_4$ to the four rows in
\eqref{eq:matrix}, and append $(0,0,2,1,2g)$ as the fifth row.
Subtracting the first row from the fifth leaves only the last entry
$2g-L_2=2(\sqrt3-1)$. The remaining determinant is that of $V$,
of absolute value one. Evaluation of a real cohomology on these
cycles therefore forces all coefficients, including the constant,
to vanish.
\end{proof}

The preceding theorems are statements about real billiard orbits, not
merely formal integer records. They also exclude a circle-valued
coboundary with a continuous transfer function on the regular return
branches: an almost-everywhere identity then extends to each regular
periodic point by continuity and full support of the section measure.
A transfer function assumed only measurable requires a separate
regularity theorem. Similarly, a covariance conclusion requires a
proved variance--coboundary criterion in the relevant function space.
If that criterion and a continuous covariance family are established,
Corollary~\ref{cor:real-period-rank} gives positive definiteness and,
on $I$, a uniform positive lower eigenvalue. Existence and continuity
of that covariance are not consequences of the determinant alone.
Compact periodic separation is not a quantitative operator contraction
or a returned temporal nonintegrability estimate.
